% Part: lambda-calculus
% Chapter: introduction
% Section: syntax

\documentclass[../../../include/open-logic-section]{subfiles}

\begin{document}

\olfileid{lam}{int}{syn}
\olsection{लॅम्डा कलनाची विन्यासमीमांसा}

सुरुवातीला $x$, $y$, $z$,~\dots अशी चरांची मालिका आणि $a$, $b$,
$c$,~\dots अशी काही स्थिरांक चिन्हे घेतली जातात. पदांचा संच पुढीलप्रमाणे
विगमनाधारित रीतीने परिभाषित केला जातो:
\begin{enumerate}
\item प्रत्येक चर हे पद आहे.
\item प्रत्येक स्थिरांक हे पद आहे.
\item $M$ आणि $N$ ही पदे असतील, तर $(MN)$ हेही पद आहे.
\item $M$ हे पद आणि $x$ हे चर असेल, तर $(\lambd[x][M])$ हे
  पद आहे.
\end{enumerate}
$(MN)$ या रूपातील पदांना \emph{उपयोजने}, तर $(\lambd[x][M])$
या रूपातील पदांना \emph{अमूर्तीकरणे} म्हणतात.

एकही स्थिरांक नसलेल्या प्रणालीला \emph{शुद्ध} लॅम्डा कलन म्हणतात.
आपण मुख्यतः शुद्ध $\lambd$-कलनात काम करणार असल्याने सर्व लहान इंग्रजी
अक्षरे चरे दर्शवतील. $\lambd$-कलनातील पदे दर्शवण्यासाठी आपण मोठी
इंग्रजी अक्षरे ($M$, $N$ इत्यादी) वापरू.

आपण संकेतनाच्या काही रूढी पाळू:

\begin{conv}
\begin{enumerate}
\item कंस वगळले असता उपयोजन डावीकडून उजवीकडे होते. उदाहरणार्थ,
  $M$, $N$, $P$ आणि $Q$ ही पदे असतील, तर $MNPQ$ हे
  $(((MN)P)Q)$ चे संक्षिप्त रूप आहे.
\item पुन्हा, कंस वगळले असता लॅम्डा अमूर्तीकरणाला शक्य तितकी
  व्यापक व्याप्ती द्यायची. उदाहरणार्थ, $\lambd[x][MNP]$ याचा
  अर्थ $(\lambd[x][((MN)P)])$ असा घ्यायचा.
\item एका लॅम्डाने अनेक चरे अमूर्त करता येतात. उदाहरणार्थ,
  $\lambd[xyz][M]$ हे
  $\lambd[x][\lambd[y][\lambd[z][M]]]$ चे संक्षिप्त रूप आहे.
\end{enumerate}
\end{conv}

उदाहरणार्थ,
\[
\lambd[xy][xxyx \lambd[z][xz]]
\]
हे पुढीलचे संक्षिप्त रूप आहे:
\[
\lambd[x][\lambd[y][((((xx)y)x)(\lambd[z][(xz)]))]].
\]
या रूढी पाठ करा. सुरुवातीला त्या तुमची डोकेदुखी वाढवतील;
पण हळूहळू त्यांची सवय होईल आणि मग कंसांच्या पसाऱ्यापेक्षा
त्या कमी त्रासदायक वाटतील.

%READERNOTE{SOL6-A102: फक्त कोणत्याही लॅम्डाच्या व्याप्तीत येणे हा बद्ध वापराचा निकष नाही. त्याच नावाचे संबंधित अमूर्तीकरण त्या वापराला बांधते. नामांतर करताना मुक्त वापर आणि इतर बंधन-संबंध जपले आहेत.}
इतर चरांचा बंधन-संबंध न बदलता ज्या दोन पदांत फक्त
बद्ध चरांची नावे बदलली आहेत, त्यांना
$\alpha$-सममूल्य म्हणतात; उदाहरणार्थ, $\lambd[x][x]$ आणि
$\lambd[y][y]$. त्यांना “तेच” पद मानणे सोयीचे ठरेल; म्हणजे
$M$ आणि $N$ ही पदे एकच आहेत असे म्हणताना “बद्ध चरांच्या
नामांतराचा फरक सोडून” ती एकच आहेत, असाही अर्थ अभिप्रेत असेल.
हे नामांतर करताना मुक्त चर पकडले जाणार नाहीत याची काळजी घेतली जाते.
एखाद्या चराचा वापर त्याच नावाच्या चराला अमूर्त करणाऱ्या
$\lambd$ च्या व्याप्तीत असेल, तर तो वापर “बद्ध” असतो;
असे बंधन नसलेला वापर “मुक्त” असतो. एकाच वापराला अनेक
अमूर्तीकरणे लागू असतील, तर सर्वांत आतले त्याच नावाचे अमूर्तीकरण
त्याला बांधते. आधीच्या उदाहरणात मुक्त चरे नाहीत; पण
\[
(\lambd[z][yz])x
\]
यात $y$ आणि $x$ मुक्त आहेत, तर $z$ बद्ध आहे.

\end{document}
